% BANKS-SOURCE: pdf=147 print=137 kind=chapter id=chapter09-opening
\BanksChapter{9}{Renormalization and effective field theory}

% BANKS-SOURCE: pdf=147 print=137 kind=prose
We are finally ready to face up to the fact that most of the expressions we
write in the Feynman-diagram solution of quantum field theory are nonsensical,
which is to say infinite. The purpose of this introductory section is to
demonstrate the simple and general counting rules which prove that this is so,
and to give the reader a conceptual orientation to the general theory of
renormalization, which we will discuss in the rest of this chapter.

We will begin with the conceptual. The formalism of quantum field theory makes
statements about arbitrarily short distances in space-time and arbitrarily high
energies. At a given moment in history experiment will never probe arbitrarily
short distances or high energies. We should expect that the formalism will have
to change to fit the data as we uncover more empirical facts about short
distances. This has happened before. Hydrodynamics is a field theory, but we
know that at short enough distance scales it is not a good description of
water. A better one is given (in principle) by the Schrödinger equation for
hydrogen and oxygen nuclei interacting with electrons, and that in turn is
superseded by the standard model of particle physics as we go to even shorter
distances and higher energies. Yet only a fool would imagine that one should
try to understand the properties of waves in the ocean in terms of
Feynman-diagram calculations in the standard model, even if the latter
understanding is possible ``in principle.''

This simple parable illustrates the idea of an \emph{effective field theory}
(EFT). EFT is a description of phenomena at wavelengths longer than some
effective cut-off scale $l_c$ and energies below some energy cut-off $E_c$. In
a theory that is (even approximately) relativistically invariant, the cut-offs
are related by $l_c\sim E_c^{-1}$. In principle, any EFT is a quantum theory,
but it may be that, as in the case of hydrodynamics, the classical
approximation is valid throughout the range of validity of the EFT itself.
The basic idea of an effective field theory is that physics in a certain
energy regime, and a certain resolution for length scales, should be described
most simply by a set of effective degrees of freedom appropriate to that scale
and depend on any underlying, more ``fundamental'' description only through a
small set of parameters governing the dynamics of these degrees of freedom.

The basic idea of renormalization is to parametrize the effect of all possible
modifications of the theory at higher energy and shorter distance in terms of
all possible interactions between the effective degrees of freedom. The
non-trivial part of renormalization theory is the mathematical demonstration
of universality of long-distance,

% BANKS-SOURCE: pdf=148 print=138 kind=prose
low-energy physics. That is, one shows that, starting from any set of
interactions at the cut-off scale, the low-energy, long-distance physics
depends on only a few \emph{relevant} parameters. The mechanics of this
demonstration involves a procedure for reducing the cut-off scale and finding
new Lagrangians that reproduce all of the same low-energy physics. Continuous
changes in the cut-off lead to a flow on the space of all possible
interactions
\[
  \frac{\dd g^i}{\dd t}=\beta^i(\{g^j(t)\}),
\]
where $t$ parametrizes an infinitesimal rescaling of the cut-off. This is
called the \emph{renormalization group (RG) flow}, despite the fact that it is
irreversible, and therefore only a semigroup. For a large class of models, the
RG can be shown to be a gradient flow, so that the general asymptotic behavior
is a \emph{fixed point}. The relevant parameters describe the manifold of
unstable directions of the fixed point. They are the parameters which have to
be tuned in order to obtain low-energy physics that is independent of the
details of the physics at the cut-off scale, and all dependence on the cut-off
is encoded in the value one chooses for these parameters.

\BanksImplicitHook{I-CH09-001}

All possible cut-off-independent infrared behaviors of a certain set of
degrees of freedom thus fall into universality classes corresponding to the
possible fixed points of their RG flow. The fixed-point Lagrangians obviously
have to be invariant under space-time scaling transformations. For most
interesting models they are actually invariant under the larger group of
space-time conformal transformations, (SO(2,$D$) for a $D$-dimensional
Minkowski space-time). Such conformal field theories (CFTs) can often be
constructed without recourse to Lagrangians and Feynman diagrams. Indeed,
Lagrangians and Feynman diagrams are concepts appropriate only to the
\emph{Gaussian CFTs}, a fancy name for free massless field theories. This leads
to a possible mathematical definition of all possible QFTs as perturbations of
CFTs by turning on non-zero values of relevant parameters. But we are getting
ahead of ourselves, and will return to these issues when we understand a
little more of the technicalities of the RG.

\BanksImplicitHook{I-CH09-002}

The basic idea behind the derivation of the RG flow equations is the concept
of \emph{integrating out degrees of freedom}. We can illustrate this in a
simple example from classical physics by imagining two harmonic oscillators
with frequencies $\Omega$ and $\omega\ll\Omega$, coupled by an anharmonic
interaction of the form $gXx^2$. We can try to solve the equations of motion
for this system by first solving for $X$ and writing equations of motion for
$x$ that take into account the evolution of $X$. We've cleverly chosen the
interaction so that the first step can be done in closed form. We find the
equations
\[
  \ddot x(t)+\omega^2x(t)=2gx(t)\int\!\dd s\,[G(t-s)x^2(s)],
\]
where the boundary conditions on $X(t)$ are encoded in the choice of Green
function $G$:
\[
  (\partial_t^2+\Omega^2)G(t-s)=\delta(t-s).
\]
In QFT we will always assume that the high-frequency degrees of freedom are
in their ground state, so the Green function is the one defined by Feynman,
which is the analytic
